\documentclass{article}
\usepackage{mdposter}
\posttitle{Slow particles and weak collisions in gyrokinetics}
\postfield{Physics, plasma physics}

\begin{document}
\begin{mdposter}

\posterbox{name=c1, column=1, between=top and bottom}{%
\begin{mdcol}
\mdsec{The question}
Gyrokinetic models describe small-scale turbulence in magnetized plasmas. In a linear electrostatic
gyrokinetic model, can collisionless transport, or weak collisions at a rate $\nu$, make the states a
computation does not resolve decay at one uniform rate?
\mdgap
\mdsec{How we are approaching it}
\begin{itemize}
\item Look for states the electric field hardly sees. Slow particles give whole infinite-dimensional subspaces whose field response vanishes uniformly.
\item Build approximate eigenvectors on the imaginary axis from differentiated wave packets. Under a nondegeneracy condition they concentrate near the passing boundary.
\item Bound their residuals for two collision models: Maxwellian velocity diffusion and a one-species linearized Coulomb operator.
\end{itemize}
\mdgap
\mdsec{What we have so far}
\begin{progress}
Quasineutral coupling does not remove the obstruction created by slow particles when the quadratic energy is coercive.
\end{progress}
\begin{progress}
If one species also has the right low-energy drift parity and an even gyroaverage, field feedback and finitely many bounded observations cannot give a decay rate uniform in $\nu$.
\end{progress}
\end{mdcol}
}

\posterbox{name=c2, column=2, between=top and bottom}{%
\begin{mdcol}
\mdsec{Spectrum of the undriven model}
\includegraphics[width=\linewidth]{boundary-spectra.pdf}
\figcap{Figure 1.}{With zero inflow at the boundary the undriven spectrum is the closed left half-plane; with periodic boundaries it is the imaginary axis.}
\mdgap
\mdsec{Residual bound}
\includegraphics[width=\linewidth]{boundary-residual.pdf}
\figcap{Figure 2.}{The curves show the residual bound stated below as the collision rate falls, with $C = 1$, at three frequencies.}
\begin{progress}
At bounded frequencies the wave-packet residuals are at most $C\bigl(\nu^{1/3} + \nu^{1/5}|\omega|^{2/5}\bigr)$, so they vanish as $\nu \to 0$.
\end{progress}
\end{mdcol}
}

\posterbox{name=c3, column=3, between=top and bottom}{%
\begin{mdcol}
\mdsec{In Mantis}
\mantis{gyrokinetic-boundary-spectra.jpg}{1,575}{Smooth drift alone does not guarantee exact zero attainment.}{%
The record is a rejected inference: a smooth-drift example comes arbitrarily close to zero frequency without reaching it.\alsoonrecord A revision rejected fixed-excess surjectivity, because nonzero frequencies stay bounded below by $2\pi/T^{*}(c^{*})$.}
\mdgap
\mdsec{Open questions}
\begin{itemize}
\item These results are for the linear models shown and the full distribution norm. Which of them survive in weaker norms?
\item Concentration near the passing boundary needs a nondegeneracy condition. Can it be dropped?
\end{itemize}
\end{mdcol}
}

\end{mdposter}
\end{document}
